\documentclass[11pt, a4paper]{article}

\usepackage[a4paper, top=2.0cm, bottom=2.0cm, left=1.8cm, right=1.8cm]{geometry}
\usepackage{fontspec}
\usepackage[english]{babel}
\setmainfont{Noto Sans}

\usepackage{amsmath}
\usepackage{amssymb}
\usepackage{amsthm}
\usepackage{booktabs}
\usepackage{graphicx}
\usepackage{url}
\usepackage{hyperref}

\emergencystretch=1.5em
\newtheorem{theorem}{Theorem}
\newtheorem{proposition}{Proposition}
\urlstyle{tt}

\hypersetup{
    colorlinks=true,
    linkcolor=blue,
    urlcolor=blue,
    citecolor=blue,
    pdftitle={A Causality-Preserving, Stereographically Compactified Worldline Representation for Retarded Lienard-Wiechert Electrodynamics},
    pdfauthor={Tim Huff}
}

\title{\textbf{A Causality-Preserving, Stereographically Compactified Worldline Representation for Retarded Li\'enard--Wiechert Electrodynamics}}
\author{\textbf{Tim Huff} \\ \small Independent Researcher (Creve Coeur, MO, USA) \\ \small Software: \url{https://github.com/timsdsdio/hopf-hough} (MIT License)}
\date{Version 6 --- October 2026}

\begin{document}

\maketitle

\begin{abstract}
\noindent We describe and verify a method, released as the open-source tool \texttt{hopf-hough}, for evaluating the exact retarded Li\'enard--Wiechert (LW) field of moving point charges at every pixel of an observation plane on the GPU. Each pixel solves, for every source, the light-cone equation $f(R) = R - \|\mathbf{x} - \mathbf{z}(R)\| = 0$ on a compact per-source history, then evaluates the LW field from the interpolated history. This is a gather-side formulation, efficient when sources are few and observation points many. Three ingredients make it accurate in single-precision WebGL2: (i) the history is resampled onto shells of retarded radius $R = c(t - t_{\text{ret}})$ placed by the stereographic map $R = L\tan(\theta/2)$; (ii) cubic Hermite segments are limited by a closed-form certificate on the Bernstein coefficients of $\|\mathbf{z}'(R)\|^2$, which guarantees $\|\mathbf{z}'\| \le \beta_{\max} < 1$ and hence a unique retarded root per observer; (iii) the root function is evaluated in double-float arithmetic on high/low knot positions, which removes a single-precision accuracy floor whose cause we identify. Against the exact field (no softening), with $M_T = 512$ shells, the relative $L_2$ error across a synchrotron pulse is $0.0095\%$ at $\gamma = 22.4$ and $0.069\%$ at $\gamma = 70.7$, compared with $0.22\%$ and $7.2\%$ for the algorithm of version 5 of this preprint; at $\gamma = 224$ it is $1.3\%$ ($M_T = 512$) and $0.45\%$ ($M_T = 2048$). The double-precision reference agrees with the independent code PyCharge to $\sim\!10^{-11}$. A validity indicator computed in the shader marks observers outside the stored history or where it is under-resolved; in far-field tests every unmarked observer is within $0.61\%$ of the exact field. No new physics is claimed.

\vspace{0.15cm}
\noindent \textbf{Keywords:} Li\'enard--Wiechert fields, retarded time, causality-preserving interpolation, Bernstein polynomials, double-float arithmetic, stereographic compactification, synchrotron radiation, WebGL2
\end{abstract}

\section{Introduction and relationship to prior work}

The field of a moving point charge is given in closed form by the Li\'enard--Wiechert expressions evaluated at the retarded time \cite{jackson1999}. Computing it numerically at many observation points is routine in accelerator physics and in teaching, and existing tools take three routes.
\begin{enumerate}
    \item \textbf{Grid solvers (FDTD/PIC)} such as WarpX \cite{vay2018} evolve $(\mathbf{E}, \mathbf{B})$ on a mesh. They handle self-consistent many-particle dynamics but suffer numerical dispersion and need spatial resolution of order $R_{\text{orb}}\gamma^{-3}$ to resolve a synchrotron pulse.
    \item \textbf{Wavefront and retarded-time LW codes.} Shintake's real-time animation \cite{shintake2003} and LANL's CoSyR \cite{huang2021} construct fields from wavefronts emitted along stored trajectories. PyCharge \cite{pycharge2022} and the accompanying visualization work \cite{filipovich2021} compute exact time-domain fields from arbitrary 3D trajectories by solving the retarded-time equation per observation point, in batch Python.
    \item \textbf{Coordinate compactification} such as $R = L\tan(\theta/2)$ is standard in numerical relativity for mapping unbounded domains to finite grids \cite{zenginoglu2008}.
\end{enumerate}

\paragraph{Gather versus scatter.} A grid solver stores the field per cell and lets sources scatter into it. The method here stores a compact history per source and lets each observation point gather: it asks each source which event on its worldline lies on the observer's past light cone. The same split between per-observer ``retarded-time'' and per-source ``source-time-dominant'' formulations is long established for moving-source noise prediction in aeroacoustics \cite{brentner1996, brentner1997}, and per-point retarded-time solves are what PyCharge does. What this work contributes is an engineering combination: a fixed-cost, single-precision GPU kernel with a guaranteed unique root, a compact radial history, and measured, reproducible accuracy. The gather formulation is efficient when the number of sources is modest and observation points are dense; grid solvers remain the right tool for self-consistent plasmas with $\sim\!10^9$ macroparticles.

\subsection*{Contributions}
\begin{itemize}
    \item A closed-form, branch-light causality certificate for cubic Hermite worldline segments based on the Bernstein coefficients of $\|\mathbf{z}'\|^2$ (Proposition~\ref{prop:sq}). It is never stricter than the control-vector limiter of version 5 (Theorem~\ref{thm:ctrl}), which we show to be about $4\times$ conservative and harmful on smooth arcs.
    \item A diagnosis of the single-precision accuracy floor (Section~\ref{sec:precision}): it comes from one subtraction in the root function and from storing knot positions in absolute coordinates, not from the $\kappa^{-3}$ factor or Newton conditioning as version 5 stated. A double-float root function on high/low knot positions removes it.
    \item An open, single-file browser tool and TypeScript library with a built-in verification suite, an independent cross-check against PyCharge, and a shader-side validity indicator.
\end{itemize}

\section{Formulation}

\subsection{Retarded Li\'enard--Wiechert field and a cancellation-free kernel}
For a point charge $q$ with trajectory $\mathbf{z}(t_{\text{ret}})$, $\boldsymbol{\beta} = \dot{\mathbf{z}}/c$ and $\dot{\boldsymbol{\beta}} = \ddot{\mathbf{z}}/c$, let $R = c(t - t_{\text{ret}})$ be the retarded radius at observation time $t$, so that $d\mathbf{z}/dR = -\boldsymbol{\beta}$. The retarded radius $R_{\text{ret}}$ for an observer at $\mathbf{x}$ solves
\begin{equation}
f(R) = R - \|\mathbf{x} - \mathbf{z}(R)\| = 0, \qquad \hat{\mathbf{n}} = \frac{\mathbf{x} - \mathbf{z}(R_{\text{ret}})}{\|\mathbf{x} - \mathbf{z}(R_{\text{ret}})\|},
\label{eq:rootfn}
\end{equation}
and the field is \cite{jackson1999}
\begin{equation}
\mathbf{E}(\mathbf{x}, t) = \frac{q(1 - \beta^2)}{\kappa^3 R^2}(\hat{\mathbf{n}} - \boldsymbol{\beta}) + \frac{q}{c\,\kappa^3 R}\hat{\mathbf{n}} \times \left[(\hat{\mathbf{n}} - \boldsymbol{\beta}) \times \dot{\boldsymbol{\beta}}\right], \qquad \kappa = 1 - \hat{\mathbf{n}}\cdot\boldsymbol{\beta}, \qquad \mathbf{B} = \hat{\mathbf{n}} \times \mathbf{E}.
\label{eq:lw}
\end{equation}
Charges move in the $x$--$y$ plane and the field is evaluated in that plane, where $\mathbf{E}$ is in-plane and $\mathbf{B}$ is along $z$. The tool offers an optional Plummer softening $R^2 \to R^2 + \epsilon^2$ for display near charges; all accuracy figures in this paper use $\epsilon = 0$. (Version 5 compared against a reference that applied the same softening, so its ``exact'' reference was not the exact field; see Section~\ref{sec:changes}.)

To avoid cancellation in $1 - \hat{\mathbf{n}}\cdot\boldsymbol{\beta}$ and $1 - \beta^2$ as $\beta \to 1$, the history stores the proper velocity $\mathbf{u} = \gamma\boldsymbol{\beta}$, recovers $1 - \beta^2 = (1 + \|\mathbf{u}\|^2)^{-1}$, and evaluates
\begin{equation}
\kappa = \frac{1 - \beta^2}{1 + \beta} + \frac{1}{2}\beta\,\|\hat{\mathbf{n}} - \hat{\boldsymbol{\beta}}\|^2, \qquad \hat{\mathbf{n}} - \boldsymbol{\beta} = \frac{1 - \beta^2}{1 + \beta}\,\hat{\mathbf{n}} + \beta\,(\hat{\mathbf{n}} - \hat{\boldsymbol{\beta}}).
\label{eq:kappa}
\end{equation}

\subsection{Stereographic radial shells}
Each source's history is resampled onto $M_T$ cell-centred shells $\theta_m = (m + \tfrac12)\Delta\theta$, $\Delta\theta = \pi/M_T$, at
\begin{equation}
R(\theta) = L\tan\!\left(\frac{\theta}{2}\right), \qquad \Delta R(\theta) = \frac{L^2 + R^2}{2L}\,\Delta\theta.
\label{eq:stereo}
\end{equation}
Spacing is nearly uniform for $R \lesssim L$, where radiation forms, and grows as $R^2$ beyond; the outermost shell is at $L\tan(\tfrac{\pi}{2} - \tfrac{\Delta\theta}{4}) \approx 4LM_T/\pi$. The map provides reach within a fixed texture, not uniform resolution to infinity: accuracy falls once $\Delta R$ no longer resolves the motion (Section~\ref{sec:far}). Shells are sampled from an append-only ring buffer in retarded time, so the textures are a resampled view of that buffer, refreshed each step.

\subsection{Causality certificates for cubic Hermite segments}
On a segment $[R_1, R_2]$, $\Delta R = R_2 - R_1$, with knot positions $\mathbf{z}_1, \mathbf{z}_2$ and knot velocities $\boldsymbol{\beta}_1, \boldsymbol{\beta}_2$ ($\|\boldsymbol{\beta}_i\| \le \beta_{\max} < 1$), write $\mathbf{v} = (\mathbf{z}_2 - \mathbf{z}_1)/\Delta R$ and assume the chord is subluminal, $\|\mathbf{v}\| \le \beta_{\max}$. A cubic Hermite curve with endpoint $R$-derivatives blended toward the secant by $\alpha \in [0,1]$,
\begin{equation}
\mathbf{V}_0(\alpha) = (1-\alpha)\mathbf{v} - \alpha\boldsymbol{\beta}_1, \qquad \mathbf{V}_2(\alpha) = (1-\alpha)\mathbf{v} - \alpha\boldsymbol{\beta}_2, \qquad \mathbf{V}_1(\alpha) = \mathbf{v} + \alpha(2\mathbf{v} + \boldsymbol{\beta}_1 + \boldsymbol{\beta}_2),
\label{eq:ctrls}
\end{equation}
interpolates the knots for every $\alpha$, and its derivative is exactly the quadratic Bernstein form $\mathbf{z}'(s) = \sum_{k=0}^{2} B_{k,2}(s)\mathbf{V}_k(\alpha)$, $s = (R - R_1)/\Delta R$.

\begin{theorem}[Control-vector certificate, version 5]\label{thm:ctrl}
$\mathbf{V}_0$ and $\mathbf{V}_2$ lie in the ball $\|\cdot\| \le \beta_{\max}$ for all $\alpha$ by convexity. If $\alpha^*$ is the largest $\alpha \in [0,1]$ with $\|\mathbf{V}_1(\alpha)\| \le \beta_{\max}$ (a closed-form quadratic root), then $\|\mathbf{z}'(R)\| \le \beta_{\max}$ on the segment by the convex-hull property of Bernstein polynomials.
\end{theorem}

The bound is a known technique: velocity limits through Bernstein control points are standard in CNC and robot trajectory planning \cite{farouki2008, tordesillas2020}, and blending tangents toward the secant is a vector form of monotonicity-preserving Hermite interpolation \cite{fritsch1980}. Applied to retarded time it is useful but conservative: for a circular arc of turning angle $\varphi$, $\|\mathbf{V}_1\| = \beta(1 + \varphi^2/8 + O(\varphi^4))$, so the limiter fires on smooth arcs whose true speed never exceeds $\beta$. On the version-5 benchmark it fired on 14--87\% of segments that were never superluminal and caused most of the $\gamma = 70.7$, $M_T = 128$ error (Section~\ref{sec:results}).

\begin{proposition}[Squared-speed certificate]\label{prop:sq}
$\|\mathbf{z}'(s)\|^2 = \sum_{k=0}^{4} B_{k,4}(s)\,c_k$ with
\[
c_0 = \mathbf{V}_0\!\cdot\!\mathbf{V}_0,\quad c_1 = \mathbf{V}_0\!\cdot\!\mathbf{V}_1,\quad c_2 = \tfrac13\left(2\,\mathbf{V}_1\!\cdot\!\mathbf{V}_1 + \mathbf{V}_0\!\cdot\!\mathbf{V}_2\right),\quad c_3 = \mathbf{V}_1\!\cdot\!\mathbf{V}_2,\quad c_4 = \mathbf{V}_2\!\cdot\!\mathbf{V}_2 .
\]
Each $c_k(\alpha)$ is quadratic in $\alpha$, and $c_k(0) = \|\mathbf{v}\|^2 \le \beta_{\max}^2$. Let $\alpha_{\text{sq}}$ be the smallest $\alpha \in (0, 1]$ at which some $c_k(\alpha) = \beta_{\max}^2$ (or $1$ if none). Then for every $\alpha \le \alpha_{\text{sq}}$, $\|\mathbf{z}'\| \le \beta_{\max}$ on the segment, and $\alpha_{\text{sq}} \ge \alpha^*$.
\end{proposition}
\begin{proof}
The product of two quadratic Bernstein expansions is a quartic Bernstein expansion with the stated coefficients; since $B_{k,4} \ge 0$ and $\sum_k B_{k,4} = 1$, $\|\mathbf{z}'\|^2 \le \max_k c_k \le \beta_{\max}^2$. $c_0$ and $c_4$ never exceed $\beta_{\max}^2$ by Theorem~\ref{thm:ctrl}, and by continuity no $c_k$ exceeds it before its first crossing. If all $\|\mathbf{V}_k\| \le \beta_{\max}$ then by Cauchy--Schwarz every $c_k \le \beta_{\max}^2$, so no crossing occurs before $\alpha^*$.
\end{proof}

Computing $\alpha_{\text{sq}}$ costs three quadratic solves. With the speed bound on every segment, $f'(R) = 1 + \hat{\mathbf{n}}\cdot\mathbf{z}'(R) \ge 1 - \beta_{\max} > 0$, so $f$ is strictly increasing on the stored domain and each observer has exactly one retarded root there. The limited worldline is $C^0$ everywhere and $C^1$ at a knot when neither neighbouring segment is limited; where a limiter acts, the velocity direction can be discontinuous at the knot. A globally $C^1$ subluminal cubic does not always exist (for alternating chords $\pm c$ with free knot slopes it requires $c \le 2\beta/3$), whereas raising the degree restores it; we leave that for future work.

\subsection{Single-precision accuracy and its cause}\label{sec:precision}
We attributed the error of a float32 implementation by promoting one stage at a time to float64 \cite{swarm2026}. At $\gamma = 70.7$ ($M_T = 512$), rounding only the root-function arithmetic $\mathbf{x} - \mathbf{z}$, $\|\cdot\|$ and $R - \|\cdot\|$ to float32 produces 12.4\% error; storing knot positions in absolute float32 coordinates produces 2.4\%; every other stage ($\kappa$ and the LW kernel, Newton conditioning, the velocity and acceleration interpolation, shell radii) contributes at most 0.05\%. The cause is conditioning: near the pulse crest $f'(R) \approx \kappa \sim \gamma^{-2}$, so an absolute error $\delta$ in $f$ moves the retarded point by $\delta/\kappa$, and float32 positions near $|\mathbf{x}| \sim 1$ carry $\delta \sim 6\times10^{-8}$.

The fix is local. Knot positions are stored as pairs $\mathbf{z} = \mathbf{z}_{\text{hi}} + \mathbf{z}_{\text{lo}}$ of float32 values, chords are formed from both parts, and $f(R)$ is evaluated with error-free transformations (double-float arithmetic \cite{thall2006, dagraca2006}); everything else stays float32. To prevent compilers from simplifying the error-free transformations, intermediate results are multiplied by a uniform equal to 1.

Two further float32 effects appear above $\gamma \approx 100$ and were corrected in this version. (i) The chord-speed clamp that guards against superluminal input compared float32 differences of adjacent knots; at $M_T = 2048$ their relative rounding ($\sim\!3\times10^{-5}$) exceeds $1 - \beta$, so physical knots were wrongly moved. The clamp now uses the full high/low positions in double precision. (ii) The Newton derivative was floored at $5\times10^{-5}$, larger than the true $f'$ at $\gamma = 224$; the floor is now $\tfrac12(1 - \beta_{\max})$, which the certificate guarantees is below $f'$.

\paragraph{How much accuracy is available.} At $\gamma = 70.7$ the pulse crossed by the benchmark sweep is about 15 float32 ulps wide in observer position, and moving an observer by one ulp changes $\mathbf{E}$ by up to 71\% of its rms value. Rounding the observer coordinates alone shifts the exact field by 3.2\% (Table~\ref{tab:gamma}). A float32-addressed method can therefore only be judged against the exact field at the rounded point, which is the metric used here, or by placing observers with double-float coordinates, which the tool supports and which is required above $\gamma \approx 100$.

\subsection{Root search}
Because $f$ is monotone at the stored knots (subluminal chords), the solver first bisects over knot indices, evaluating $f$ at knots only (one texture fetch each), and then runs safeguarded Newton inside the bracketing segment on the certified cubic, iterating to convergence (at most 12 steps, with bisection fallback). The segment index and the fraction within it are kept as separate values, so float32 resolution of a combined coordinate never limits the retarded point.

\subsection{Validity indicator}
For each source the shader reports $q = \max(\Delta R\,\gamma\,|\dot{\beta}_\perp|/\beta^2,\; 6\,\|\Delta\mathbf{u}\|/(\bar{u} + 0.1))$ at the retarded point: shell spacing times $\gamma$ over the local radius of curvature, and the relative change of proper velocity across the segment. It also marks retarded points beyond the stored history ($q = -1$), inside the first shell, and on shells damped by the optional Nyquist taper. The app hatches pixels with $q > 1.5$ or $q < 0$. The thresholds are heuristic; Section~\ref{sec:far} measures how well they predict the error.

\section{Verification protocol}\label{sec:protocol}
All numbers come from scripts in the repository (\texttt{tests/gpu/tables.mjs}) and can be regenerated in a browser with the in-app suite.
\begin{enumerate}
    \item \textbf{Reference.} The exact field is computed in float64 directly on the analytic trajectory (32-step bracketed bisection on $t_{\text{ret}}$, then up to 12 safeguarded Newton steps) with $\epsilon = 0$. With softening off it agrees with PyCharge 2.0.1 \cite{pycharge2022}, which obtains velocity and acceleration by automatic differentiation of the position, to at most $1.3\times10^{-11}$ relative error for circular orbits ($\beta = 0.5, 0.9, 0.99$, including pulse sweeps), harmonic oscillation and straight-line motion.
    \item \textbf{Benchmark.} A charge on a circular orbit of radius $0.85$ ($c = 2.35$). A ring buffer ($\Delta t = 0.00675$) is seeded and advanced 16 steps; the shells are populated from it. Observers: 256 points at fixed distance from the orbit centre crossing the pulse emitted tangentially at retarded radius $1.15$ (half-width $4\gamma^{-3}$ rad), and a 64-point ring at radius $1.15$.
    \item \textbf{Readback.} Fields are read back from an \texttt{RGBA32F} framebuffer (\texttt{gl.readPixels}). The measurements here use WebGL2 with the SwiftShader software renderer, which implements IEEE float32; real GPUs should be checked with the in-app suite.
    \item \textbf{Regression.} The single-file solver deposited with version 5 is kept in the repository. The TypeScript port reproduces all 154 values of its verification suite exactly in the same browser engine; the version-5 column in Table~\ref{tab:gamma} is that algorithm run on the same history but against the unsoftened reference.
\end{enumerate}

\section{Results}\label{sec:results}

\subsection{Accuracy across Lorentz factors}
\begin{table}[htbp]
\centering
\resizebox{\textwidth}{!}{%
\begin{tabular}{@{}rrrrrrrr@{}}
\toprule
$\beta$ & $\gamma$ & $M_T$ & \textbf{this version} & \textbf{this version, df64 observers} & version-5 algorithm & ring (this version) & observer rounding alone \\
\midrule
0.8     & 1.67  & 128  & 0.0067\%            & 0.0067\%            & 0.0067\% & 0.0030\%            & $2.2\times10^{-5}$\% \\
0.8     & 1.67  & 512  & $4.5\times10^{-4}$\% & $4.4\times10^{-4}$\% & 0.0010\% & $1.5\times10^{-4}$\% & \\
0.8     & 1.67  & 2048 & $4.5\times10^{-5}$\% & $4.6\times10^{-5}$\% & 0.0034\% & $2.8\times10^{-5}$\% & \\
0.95    & 3.20  & 128  & 0.017\%             & 0.017\%             & 0.017\%  & 0.12\%              & $3.1\times10^{-4}$\% \\
0.95    & 3.20  & 512  & 0.0010\%            & 0.0010\%            & 0.0019\% & 0.0063\%            & \\
0.95    & 3.20  & 2048 & $1.5\times10^{-4}$\% & $1.3\times10^{-4}$\% & 0.0075\% & $3.8\times10^{-4}$\% & \\
0.99    & 7.09  & 128  & 0.038\%             & 0.038\%             & 0.039\%  & 0.012\%             & 0.0034\% \\
0.99    & 7.09  & 512  & 0.0023\%            & 0.0023\%            & 0.0097\% & $8.6\times10^{-4}$\% & \\
0.99    & 7.09  & 2048 & $4.3\times10^{-4}$\% & $6.0\times10^{-4}$\% & 0.015\%  & $8.1\times10^{-5}$\% & \\
0.999   & 22.4  & 128  & 0.13\%              & 0.13\%              & 0.26\%   & 0.013\%             & 0.11\% \\
0.999   & 22.4  & 512  & \textbf{0.0095\%}   & 0.0096\%            & 0.22\%   & 0.0016\%            & \\
0.999   & 22.4  & 2048 & 0.0067\%            & 0.0055\%            & 0.24\%   & $3.3\times10^{-4}$\% & \\
0.9999  & 70.7  & 128  & 0.37\%              & 0.43\%              & 48\%     & 0.045\%             & 3.2\% \\
0.9999  & 70.7  & 512  & \textbf{0.069\%}    & 0.069\%             & 7.2\%    & 0.0027\%            & \\
0.9999  & 70.7  & 2048 & 0.051\%             & 0.043\%             & 5.6\%    & $1.7\times10^{-4}$\% & \\
0.99999 & 224   & 128  & ---                 & 7.0\%               & ---      & 0.041\%             & $>100$\% \\
0.99999 & 224   & 512  & ---                 & \textbf{1.3\%}      & ---      & $9.6\times10^{-4}$\% & \\
0.99999 & 224   & 2048 & ---                 & \textbf{0.45\%}     & ---      & $1.2\times10^{-4}$\% & \\
\bottomrule
\end{tabular}}
\caption{Relative $L_2$ error of GPU readback against the exact LW field (no softening). Sweep: 256 observers crossing the synchrotron pulse; ring: 64 observers at radius 1.15. ``This version'' rounds observers to float32 and compares with the exact field at the rounded point; ``df64 observers'' places them with double-float coordinates. At $\gamma = 224$ float32 observer coordinates cannot resolve the pulse (last column), so only the df64 column is meaningful. The version-5 algorithm uses a float32 root function, the Theorem~\ref{thm:ctrl} limiter and bisection on the spline; it was not designed for $\gamma > 100$.}
\label{tab:gamma}
\end{table}

Table~\ref{tab:gamma} shows a roughly $100\times$ reduction in error at $\gamma = 22$--$71$ relative to the version-5 algorithm, with the largest change at $M_T = 128$, $\gamma = 70.7$ (48\% to 0.37\%), where the version-5 limiter was active across the pulse. At $\gamma \le 7$ both algorithms are accurate; the remaining error there is interpolation of the history and converges with $M_T$.

\subsection{Radial grid at equal memory}
\begin{table}[htbp]
\centering
\begin{tabular}{@{}lrrrr@{}}
\toprule
128 shells, $R \in [0, 220]$ & sweep & ring & spacing at $R = 1.15$ & first spacing \\
\midrule
Stereographic $L\tan(\theta/2)$, $L = 1.35$ & \textbf{0.038\%} & 0.012\% & 0.029 & 0.017 \\
Logarithmic $r_0(e^{\alpha u} - 1)$         & 0.33\%  & 0.080\% & 0.081 & 0.0035 \\
Quadratic $R_{\max} u^2$                    & 1.6\%   & 0.69\%  & 0.27  & 0.027 \\
Uniform                                     & $>100$\% & $>100$\% & 1.72 & 1.72 \\
\bottomrule
\end{tabular}
\caption{Same memory and reach, $\beta = 0.99$, prefilter off for every grid (version 5 applied it only to the stereographic grid, and its quadratic-grid entry of 0.21\% does not reproduce). The stereographic map's advantage here follows from Eq.~\eqref{eq:stereo}: near-uniform spacing across the region $R \lesssim L$ that contains the orbit and the observation radius, where the tested logarithmic grid concentrates shells near $R = 0$. Other parameter choices for the competing grids may narrow the gap.}
\label{tab:grids}
\end{table}

\subsection{Far field and the validity indicator}\label{sec:far}
\begin{table}[htbp]
\centering
\begin{tabular}{@{}rrrrrrr@{}}
\toprule
$M_T$ (reach) & $R = 1.15$ & $R = 3$ & $R = 10$ & $R = 30$ & $R = 100$ & $R = 200$ \\
\midrule
128 (220)   & 0.0064\% (0) & 0.61\% (0)  & 2.9\% (64)  & $>100$\% (64) & $>100$\% (64) & $>100$\% (64) \\
512 (880)   & 0.0015\% (0) & 0.011\% (0) & 0.20\% (26) & 31\% (64)     & $>100$\% (64) & $>100$\% (64) \\
2048 (3520) & $9.2\times10^{-5}$\% (0) & 0.0019\% (0) & 0.010\% (0) & 0.18\% (64) & $>100$\% (64) & $>100$\% (64) \\
\bottomrule
\end{tabular}
\caption{Circular orbit, $\beta = 0.9$: relative $L_2$ error on rings of radius $R$ around the orbit centre, and in parentheses the number of the 64 observers marked by the validity indicator. Across all far-field cases run ($\beta = 0.5$ and $0.9$; $M_T = 128$, 512, 2048; taper on and off), the largest error over unmarked observers is 0.61\%.}
\label{tab:far}
\end{table}

Accuracy falls with distance once the shell spacing, which grows as $R^2$, no longer resolves the orbital motion; the stored history reaches far beyond the radius at which it stays accurate. The indicator is conservative (it marks some observers whose error is below 1\%) but in these tests it never left a large error unmarked. The optional Nyquist taper of version 5 damps the radiation term on coarse shells; it is off by default, and the shells it modifies are marked.

\subsection{When the causality limiter matters}
On smooth histories resampled onto the shells, the unlimited cubic rarely exceeds $c$. In the synchrotron benchmark its peak speed equals the knot speed. For hyperbolic motion (constant proper acceleration over 100 time units, $\gamma \approx 11$ at the observation time) it reaches $\|\mathbf{z}'\| = 1.057$, but only on the outermost shell ($R \in [73, 220]$, $M_T = 128$). Observers placed so that their retarded points fall in that segment see 21--73\% error with the squared-speed certificate, 49--96\% with Theorem~\ref{thm:ctrl} and 68--91\% with no limiter, and the indicator marks all of them. The certificate's role is therefore a guarantee, one root per observer so the search cannot take a wrong branch, at no accuracy cost; it does not make an under-resolved history accurate.

\subsection{Cost}
In software rendering, the production kernel (double-float root function, squared-speed certificate, knot search with converged Newton) and the version-5 kernel take the same time to within the $\pm10\%$ run-to-run noise (640~$\times$~400 pixels, 16 sources, $M_T = 512$). Timings on GPU hardware are not yet reported; the in-app benchmark measures them.

\section{Limitations and scope}
\begin{enumerate}
    \item \textbf{Two dimensions and prescribed motion.} Charges move in a plane and do not respond to each other's fields; there is no radiation reaction.
    \item \textbf{History, not far-field storage.} The shell textures are a resampled view of an append-only ring buffer in retarded time, refreshed every step. Coarsening old history removes the radiation pulse, which needs full-resolution acceleration. Handing off to a stored far field on null slices $(\tau, \hat{\mathbf{n}})$ would only be valid beyond $R \gtrsim 175\,\gamma\rho/\varepsilon$ for relative error $\varepsilon$ ($\rho$ the radius of curvature) and would cost more memory than the worldline, because the far field undoes the $\gamma^3$ compression of the emission time; the worldline is close to the most compact representation. (Version 5 proposed such a hand-off at $R \approx R_{\text{Nyq}}$ and a Hopf-bundle representation of the polarization; both are withdrawn. A transverse field over the sphere needs no bundle when stored as a Cartesian vector orthogonal to $\hat{\mathbf{n}}$, and this pattern is the established Cauchy-characteristic extraction \cite{bishop1997}.)
    \item \textbf{Very high $\gamma$.} The method is tested to $\gamma = 224$. Above $\gamma \approx 100$ observers must be placed with double-float coordinates.
    \item \textbf{Hardware.} Accuracy was measured with an IEEE float32 software renderer. Compilers that reorder floating-point operations aggressively could weaken the double-float steps; the in-app suite checks this on the user's GPU.
    \item \textbf{Sources faster than the wave speed.} In media or acoustics, sources can outrun the wave and observers see several retarded roots. The certificate does not apply there; counting all roots per segment is a natural extension.
\end{enumerate}

\section{Changes from version 5}\label{sec:changes}
\begin{itemize}
    \item The ``exact'' reference in version 5 applied the Plummer softening ($\epsilon = 0.06$). All errors here are against the unsoftened field.
    \item The single-precision floor was attributed in version 5 to $\kappa^{-3}$ evaluation and Newton conditioning; the measured cause is the root-function subtraction and absolute knot storage (Section~\ref{sec:precision}). The 6--9\% GPU-versus-CPU discrepancy reported there is reproduced by legal reorderings of float32 operations and is not a hardware effect.
    \item The Theorem~\ref{thm:ctrl} limiter caused most of the $\gamma = 70.7$, $M_T = 128$ error; it is replaced by Proposition~\ref{prop:sq}. Version 5's Float64 convergence column ran without any limiter, so it measured a different algorithm from the GPU path.
    \item Version 5's far-field result (1.62\% at $R_{\text{obs}} = 45$) used a preset that moved at $0.44c$ with no acceleration and jumped back every 4.4 time units; it tested a Coulomb-like field, not radiation. It is replaced by Table~\ref{tab:far}.
    \item Some version-5 table entries do not reproduce with its own code: the quadratic grid (4.4\% measured vs 0.21\% printed), the $\beta = 0.999$, $M_T = 128$ cell (0.25\% vs 2.47\%) and the max~$|\Delta\mathbf{z}|$ column; three far-field cells were not computed by the deposited suite.
    \item The ensemble (``swarm'') preset double-counted the frame rotation in its acceleration (15--17\% error in $\dot{\boldsymbol{\beta}}$) and rescaled velocity without position; both are fixed.
    \item Prior work on aeroacoustic retarded-time formulations, velocity-bounded B\'ezier curves and PyCharge is now cited; the far-field hand-off and Hopf-bundle proposal are withdrawn (Section 5).
\end{itemize}

\section{Code and data availability}
The tool, library, tests and the scripts that produce every number in this paper are at \url{https://github.com/timsdsdio/hopf-hough} under the MIT License, together with the version-5 solver, the PyCharge cross-check and the raw results (\texttt{results/tables.json}). The app runs from a single HTML file with no installation.

\section*{Acknowledgments and AI assistance}
The physical idea --- replacing a global grid with per-source compactified histories, inspired by parameter-space (Hough) transforms --- and the direction of this work are the author's. AI systems contributed substantially, and all results were produced by code in the repository and checked by the author:
\begin{itemize}
    \item \textbf{Gemini (Google)}: % TODO(Tim): add the model version(s) used.
    the original WebGL2 shader, the Theorem~\ref{thm:ctrl} construction, the cancellation-free kernel, the version-5 verification suite and drafts of earlier versions.
    \item \textbf{Astra}: % TODO(Tim): say which system/product "Astra" is, and its version.
    adversarial review of earlier drafts and code (six rounds), including the pulse-crossing sweep and the matched-reach grid benchmark.
    \item \textbf{Claude (Anthropic), Claude Fable 5.1 and Claude Opus 5.5}: review of version 5; a multi-agent analysis that identified the limiter's conservatism and the squared-speed certificate, the cause of the single-precision floor and the double-float fix, the far-field and prior-art issues, and the preset defects; the TypeScript port, production shader, tests, PyCharge cross-check and documentation; and drafting of this version.
\end{itemize}

\begin{thebibliography}{99}

\bibitem{jackson1999}
J.~D. Jackson, \textit{Classical Electrodynamics}, 3rd ed. (Wiley, 1999), Ch.~14.

\bibitem{vay2018}
J.-L. Vay \textit{et al.}, ``Warp-X: A new exascale computing platform for beam--plasma simulations,'' \textit{Nucl. Instrum. Methods Phys. Res. A} \textbf{909}, 476--479 (2018).

\bibitem{shintake2003}
T. Shintake, ``Real-time animation of synchrotron radiation,'' \textit{Nucl. Instrum. Methods Phys. Res. A} \textbf{507}, 89--92 (2003).

\bibitem{huang2021}
C.-K. Huang, T.~J.~T. Kwan and B.~E. Carlsten, CoSyR: wavefront-based Li\'enard--Wiechert code for multi-dimensional synchrotron radiation, Los Alamos National Laboratory; code at \url{https://github.com/lanl/cosyr}.

\bibitem{pycharge2022}
M. Filipovich and S. Hughes, ``PyCharge: An open-source Python package for self-consistent electrodynamics simulations of Lorentz oscillators and moving point charges,'' \textit{Comput. Phys. Commun.} \textbf{274}, 108291 (2022).

\bibitem{filipovich2021}
M. Filipovich and S. Hughes, ``Space-time computation and visualization of the electromagnetic fields and potentials generated by moving point charges,'' \textit{Am. J. Phys.} \textbf{89}, 482--489 (2021).

\bibitem{zenginoglu2008}
A. Zengino\u{g}lu, ``Hyperboloidal evolution with the Einstein equations,'' \textit{Class. Quantum Grav.} \textbf{25}, 145002 (2008).

\bibitem{brentner1996}
K.~S. Brentner, ``Numerical algorithms for acoustic integrals --- the devil is in the details,'' AIAA Paper 96-1706 (1996).

\bibitem{brentner1997}
K.~S. Brentner, ``An efficient and robust method for predicting helicopter high-speed impulsive noise,'' \textit{J. Sound Vib.} \textbf{203}, 87--100 (1997).

\bibitem{farouki2008}
R.~T. Farouki, \textit{Pythagorean-Hodograph Curves: Algebra and Geometry Inseparable} (Springer, 2008).

\bibitem{tordesillas2020}
J. Tordesillas and J.~P. How, ``MINVO basis: Finding simplexes with minimum volume enclosing polynomial curves,'' arXiv:2010.10726 (2020).

\bibitem{fritsch1980}
F.~N. Fritsch and R.~E. Carlson, ``Monotone piecewise cubic interpolation,'' \textit{SIAM J. Numer. Anal.} \textbf{17}, 238--246 (1980).

\bibitem{thall2006}
A. Thall, ``Extended-precision floating-point numbers for GPU computation,'' ACM SIGGRAPH Research Posters (2006).

\bibitem{dagraca2006}
G. Da Gra\c{c}a and D. Defour, ``Implementation of float-float operators on graphics hardware,'' arXiv:cs/0603115 (2006).

\bibitem{bishop1997}
N.~T. Bishop, R. G\'omez, L. Lehner, M. Maharaj and J. Winicour, ``High-powered gravitational news,'' \textit{Phys. Rev. D} \textbf{56}, 6298 (1997).

\bibitem{swarm2026}
Analysis scripts, saved outputs and reports for the error attribution, the limiter study and the far-field estimates: \texttt{validation/analysis/} in the software repository.

\end{thebibliography}

\end{document}
